% ======================================================================
% SECTION 17: DISCUSSION
% File: sections/discussion.tex
% ======================================================================

[PROCESSING INSTRUCTION: Generate this section (approximately 2-3 pages) providing a
substantive discussion of the autonomous sponsor system. Use the following source files:

SOURCE FILES TO PROCESS:
- All sponsor/input\_files/ documents for comprehensive context
- national-platform/new\_paper/final\_paper/sections/discussion.tex (platform discussion)
- national-platform/new\_paper/final\_paper/sections/conclusion.tex (platform findings)
- patient-journey/paper/patient\_journey\_paper.tex (journey discussion)

SUBSECTION STRUCTURE:

\subsection{Paradigm Shift in Sponsor Operations}
- From human-staffed to AI-native sponsor operations
- Comparison with existing clinical trial software (EDC, CTMS, eTMF)
- What makes this different: company-shaped trial execution application
- Boundaries of autonomous operation: human sponsor-of-record role

\subsection{Regulatory Compatibility}
- How the adapted 21 CFR 312, 21 CFR 50, and ICH E6(R3) enable autonomous sponsorship
- Current regulatory gaps and needed guidance for AI-native sponsors
- International harmonization considerations (FDA, EMA, PMDA)
- The legal entity question: who is legally responsible?

\subsection{Physical AI Integration Challenges}
- Robotic system reliability and safety requirements
- Real-time decision-making under uncertainty
- Emergency handling and human override mechanisms
- Device regulatory pathway considerations (510(k), De Novo, PMA)

\subsection{Limitations and Future Work}
- Current system is a template/specification, not a deployed system
- Safety-critical decisions still require human approval gates
- Regulatory acceptance pathway is undefined
- Need for prospective validation studies
- Integration with existing pharmaceutical IT infrastructure
- Scalability beyond oncology to other therapeutic areas

PYTHON SCRIPT REQUIREMENTS: No new Python scripts for this section.
This is a narrative discussion section that synthesizes findings from
all preceding sections. Process with Claude Code Opus 4.6 (1M token context).

FORMATTING NOTES:
- This should be a thoughtful, substantive discussion, not a summary.
- Address both the promise and the challenges honestly.
- Reference specific adapted regulations and how they apply.
- Cross-reference relevant preceding sections.]

\subsection{Paradigm Shift in Sponsor Operations}

[PROCESSING INSTRUCTION: Generate 0.75 page discussing the paradigm shift. Emphasize
that this is not just another clinical trial technology tool. It is a company-shaped
trial execution application that replaces each sponsor function with an autonomous agent.
Contrast with existing tools (EDC systems like Medidata Rave, CTMS like Oracle Siebel,
eTMF systems). Explain the human sponsor-of-record wrapper.
Process with Claude Code Opus 4.6 (1M token context).]

\subsection{Regulatory Compatibility}

[PROCESSING INSTRUCTION: Generate 0.75 page on regulatory compatibility. Discuss how the
adapted regulatory frameworks (21 CFR 312 Subpart J for Physical AI, 21 CFR 50 Subpart C,
ICH E6(R3) adapted) provide the foundation. Identify gaps: no FDA guidance on AI-native
sponsors, liability questions, international harmonization needs. Reference the legal
sponsor-of-record question. Process with Claude Code Opus 4.6 (1M token context).]

\subsection{Physical AI Integration Challenges}

[PROCESSING INSTRUCTION: Generate 0.75 page on Physical AI challenges. Cover: robotic
system reliability requirements for clinical use, real-time safety decision-making,
emergency stop and human override mechanisms, and device regulatory pathways.
Reference the safety model from the repository.
Process with Claude Code Opus 4.6 (1M token context).]

\subsection{Limitations and Future Work}

[PROCESSING INSTRUCTION: Generate 0.75 page on limitations. Be honest about the current
state: this is a specification, not a deployed system. Discuss: need for human approval
gates in safety-critical decisions, undefined regulatory acceptance pathway, need for
prospective validation, IT infrastructure integration challenges, and potential extension
beyond oncology. Process with Claude Code Opus 4.6 (1M token context).]
